\scanpage{103}
\mspage{96}
% archive sheet number: 102
Mais revenons à ce cas
\[
X=\varinjlim X_i\quad\left|\begin{aligned}
&X_i\in M^{\pi}\\
&I\text{ grand devant }\pi
\end{aligned}\right.,
\]
et \unsure{introduisons} à nouveau
\[
\widetilde I=\{I'\subset I\mid I'\text{ filtrant, de cardinal }\leq\pi\},
\]
donc cette fois $\widetilde I$ est \deleted{\unsure{fort}} bien adapté à $\pi$ (filtrant, et toute partie \inserted{filtrante} de cardinal $\leq\pi$ a une borne supérieure), et si on pose
\[
X_{\widetilde i}=\varinjlim_{i\in\unsure{I'_{\widetilde i}|J'_{\widetilde i}}}X_i,
\]
on voit que ce syst. inductif
\[
(X_{\widetilde i})_{\widetilde i\in\widetilde I}
\]
est parfaitement adapté à $\pi$ \unsure{et|à} la filtr. cardinale. De plus, \deleted{\unsure{par une|pour une} application canonique
\[
I\longrightarrow\widetilde I\qquad i\longmapsto\{i\},
\]
(pas \unsure{croissant|croissante}, \unsure{qu'il est évident}), et $I$ est}
je dis que l'ens $\widetilde I_M$ est \uline{cofinal} dans $\widetilde I$. En effet,
\deleted{soit $I'\in\widetilde I$, et soit $i_0$ un majorant de $I'\deleted{\unsure{\subset|\in}}I$ dans $I$. Soit}
$\widetilde I_M$ contient le sous-ensemble $\widetilde I_0$ \unsure{formé|formée|fermé|fermée} des $I'\in\widetilde I$ qui ont un \deleted{plus} grand élément -- et $\widetilde I_0$ est bien cofinal! Maintenant, on peut appliquer,
